% !TeX root = 20261029_exp93.tex
% !LW recipe=lablog XeLaTeX
\input{_support/style.tex}
\usepackage{tabularx}
\hypersetup{pdftitle={EXP93: CIC2018 PCAP·netFound pipeline 점검과 전체 holdout 재실행 설계}}
\begin{document}
\reporttitle{2026-10-29 실행 · 2026-11-01 재검토}{EXP93: CIC2018 PCAP·netFound pipeline 점검과\\전체 holdout 재실행 설계}
\recordnote{실행 \path{20261029_192036_cic2018_pcap_expanded_59b3028b}은 120/120개 classifier job을 완료했지만, trainㆍvalidationㆍtest 모두 family$\times$split당 최대 1,024개 reservoir로 보존했다. 정상 포함 test는 6,328행에 그쳤고, 동일 입력ㆍ상이 라벨 감사도 19,725개 retained flow에 한정됐다. 따라서 이 실행의 F1 표와 충돌 수는 PCAP 기반 새 CIC2018의 난도ㆍ포화ㆍ라벨 결함 규모를 판정하는 근거로 사용하지 않는다. protocol은 train/validation만 cap하고 test의 모든 group-held-out flow를 보존하도록 수정했다.}

\section{목적과 이번 재검토}
\textbf{EXP93은 PCAP에서 새 flow vector를 구성했을 때 CIC2018의 두 문제를 다시 점검하려는 실험이다.} 첫째, encoder 입력이 같은데 family 라벨이 다른 flow가 얼마나 남는지 확인한다. 둘째, class당 n개 라벨만 사용해도 F1이 이르게 높아지는지 전체 holdout에서 확인한다. frozen netFound-v2-small의 headerㆍtiming 표현을 PCA 64차원으로 만들고, XGBoost와 TabPFN에 같은 n개 training flow를 제공하는 설계는 이 두 질문에 사용할 수 있다.

\textbf{완료된 실행의 test 보존 방식은 이 목적과 맞지 않았다.} group hash로 train/validation/test를 나눈 뒤 세 split 모두 label별 최대 1,024개만 남겼다. 그 결과 n을 바꿔도 평가 대상은 정상 포함 6,328행, 공격 family만 5,304행인 균형 reservoir였다. 이 표본은 pipeline 작동과 두 모델의 동일 입력 비교에는 쓸 수 있지만, 실제 held-out flow의 다양성ㆍ반복성ㆍ난이도와 포화 여부를 나타내지 않는다.

\section{완료된 pilot의 범위}
\textbf{pilot은 9개 CIC2018 host PCAP에서 1,812,228개 flow를 재구성하고 19,725개만 보존했다.} 확정 공격은 1,491,948개, 배경 정상은 317,603개, attempted-to-benign은 2,646개였다. raw annotation은 flow의 \texttt{annotation\_label}로 유지하고, 분류에서는 기존 비교와 맞춘 \texttt{benign}, \texttt{ddos}, \texttt{bot}, \texttt{dos}, \texttt{bf}, \texttt{inf}, \texttt{web}의 7개 family를 사용했다. Botnet Ares에는 public endpoint가 host PCAP에서 private VPC endpoint로 관측되는 capture 한정 alias를 적용했고, 그 alias와 rule은 protocol에 기록했다.

\begin{table}[H]
\centering\small
\caption{group split 뒤 원시 family flow 수와 pilot에 실제 남긴 train/test flow 수.}
\begin{tabular}{lrrrr}
\toprule
family & raw train & raw test & pilot train & pilot test\\\midrule
benign & 223,250 & 47,103 & 1,024 & 1,024\\
ddos   & 784,396 & 167,096 & 1,024 & 1,024\\
bot    & 13,972 & 2,712 & 1,024 & 1,024\\
dos    & 47,288 & 10,304 & 1,024 & 1,024\\
bf     & 65,874 & 14,379 & 1,024 & 1,024\\
inf    & 132,046 & 28,192 & 1,024 & 1,024\\
web    & 893 & 184 & 893 & 184\\
\bottomrule
\end{tabular}
\end{table}

\textbf{pilot F1은 full-holdout 성능이나 데이터 상한의 증거가 아니다.} n=2에서 정상 포함 macro-F1은 XGBoost 47.42\%, TabPFN 65.11\%였고, n=20에서는 89.72\%, 92.50\%였다. n=20에서 n=50까지 XGBoost는 4.35\%p, TabPFN은 2.86\%p 상승했다. 이 수치는 fixed reservoir에서 두 baseline의 학습곡선을 확인한 기록으로만 보존하며, full test 재실행 후의 curve와 대체한다.

\reportfigure{figures/exp93_macro_f1.pdf}{pilot reservoir test에서의 macro-F1. 이 그림은 full-holdout 재실행 전의 pipeline 확인 기록이며 CIC2018의 포화 판정에는 사용하지 않는다.}

\section{동일 입력ㆍ상이 라벨 감사의 범위}
\textbf{pilot의 충돌 수는 retained flow에만 해당한다.} frozen netFound 입력 전 signature에서 cross-label 그룹은 7개, 행은 34개였고, PCA 64차원 vector에서는 43개 그룹, 151개 행이었다. 모든 source flow를 보존하거나 전체 holdout을 평가한 값이 아니므로, 원본 CIC CSV의 동일 vectorㆍ상이 라벨 현상과 비율을 비교하거나 PCAP pipeline이 문제를 해결했다고 결론낼 수 없다.

\textbf{다음 audit은 complete group-held-out test population에서 input signature와 projected vector를 각각 집계한다.} input signature는 headerㆍtiming token을 직렬화한 hash로, PCA vector는 바이트 단위 hash로 그룹화한다. 각 cross-label 그룹에는 family 조합, 행 수, split별 수, flow ID 예시를 기록한다. 전체 test의 family별 cross-label 행 비율을 함께 제시해, 단순한 그룹 수가 아니라 해당 family에서 현상이 차지하는 규모를 확인한다.

\section{수정한 재실행 protocol}
\textbf{수정 protocol은 label budget만 training에 적용하고 test에는 적용하지 않는다.} deterministic bidirectional 5-tuple group hash로 먼저 70/15/15 split을 만든다. train과 validation은 family당 최대 1,024개 reservoir를 보존하지만 test의 labeledㆍunresolved flow는 모두 보존한다. few-shot training은 이 retained train pool에서 n=2, 3, 5, 10, 20, 50개를 seed별로 선택하며, 두 baseline은 매 조건에서 완전히 같은 training flow를 받는다.

\textbf{full holdout에서는 class support가 원시 test 분포를 따른다.} 정상 포함 평가의 최소 labeled test 행 수는 표의 raw test 합계인 269,970개다. macro-F1은 family별 F1의 동일 가중 평균으로 유지하되, 각 family F1은 그 family의 모든 held-out flow로 계산한다. GPU memory를 제한하기 위해 fitted classifier의 확률 계산만 8,192행 단위로 나누고, 예측을 원래 순서로 합쳐 같은 full test에서 F1을 계산한다. 따라서 DDoS처럼 행이 많은 family의 내부 변이와 Web처럼 184행뿐인 family의 불확실성이 F1 및 confidence interval에 그대로 반영된다.

\textbf{재실행은 기존 PCAP 파일을 재다운로드하지 않고 corpus부터 다시 구성한다.} 이전 reservoir는 버린 flow의 packet locator와 bounded token을 보존하지 않았으므로, 기존 \texttt{flows.jsonl}만 확장할 수 없다. PCAP를 다시 순회해 full test corpus를 만들고, frozen netFound encode, PCA, audit, XGBoostㆍTabPFN 120개 job 순으로 실행한다. 새로운 encoder는 동일 flow IDㆍsplitㆍlabel을 가진 별도 view를 작성할 수 있으므로 이후 ET-BERT나 NetMamba view를 추가해 같은 full holdout에서 비교할 수 있다.

\section{재실행 뒤 판단할 기준}
\textbf{새 결과는 두 기준을 분리해 해석한다.} n=20 부근의 높은 F1만으로 포화를 선언하지 않고, n=20$\rightarrow$50의 변화, full-test class F1, seed 간 변동, family별 support를 함께 본다. 동일 입력ㆍ상이 라벨은 full test의 family별 행 비율과 사례를 기준으로 판단한다. TabPFN의 few-shot 이득은 각 seed에서 같은 nㆍ같은 flow에 대한 TabPFN$-$XGBoost 차이와 full holdout macro/class F1으로 제시한다.

\textbf{이번 pilot은 재실행을 위한 corpusㆍlabelㆍencoder 검증을 제공한다.} PCAP fetch, flow reconstruction, capture-scoped Bot alias, group split, frozen-weight 검사와 paired training selection은 정상 종료했다. full-holdout 실행은 이 기반 위에서 표본화가 만든 해석 공백을 없애고, 새 PCAP representation이 CSV 입력의 두 문제를 실제로 줄이는지 검증한다.

\sourcepath{tabpfn/results/v7/exp93/20261029\_192036\_cic2018\_pcap\_expanded\_59b3028b/RUN\_META.json}
\sourcepath{tabpfn/results/v7/exp93/20261029\_192036\_cic2018\_pcap\_expanded\_59b3028b/corpus/manifest.json}
\sourcepath{tabpfn/results/v7/exp93/20261029\_192036\_cic2018\_pcap\_expanded\_59b3028b/encoders/netfound-v2-small/audit.json}
\sourcepath{docs/lablog/evidence/20261029/exp93\_summary.json}
\end{document}
